% Figure polypeps_cylinder_projector.
% Owner: none (note-only figure).
% Used by: gapped-2d-notes PP Sec. V A (Proposition 3, the cylinder form of P),
% unnumbered display.
% Source: arXiv OpenAI2026PolynomialPEPS, sections/03-patches.tex:24-49
% (Proposition 4.1, eq. (4.1) patch-cylinder-projection).
% Formula: P = sum_{j=1}^m sum_{l=1}^{d_j} |v_{jl}><v_{jl}|_{X_j} (x) id_{Lambda \ X_j},
% with X_j = X(r_j), u <= r_1 <= ... <= r_m <= 2u, sum_j d_j <= C L^c.
% Ink-to-index: invisible (skin=none) identity vertices only extend wires to the
% panel edges. The three wires, top to bottom, carry H_{X_j}, H_{X(2u) \ X_j}
% and H_{Lambda \ X(2u)}; they are named once, at the east end of the right
% panel. Left: P is supported on X(2u) (the west bracket), so its box spans the
% top two wires and the bottom wire is the identity; P is inscribed (size l).
% The left panel's west legs are free ports (no open-w wires), so the bracket
% X(2u) sits at the very ends of the two legs it names instead of across a
% longer stub; a void bracket column would push the display past one column.
% Right: the (j,l) summand consumes its input on X_j by the bra <v_{jl}| (east)
% and prepares the ket |v_{jl}> (west), in adjacent columns; the two boxes are
% separate vertices with no wire between them, and the other two wires pass
% straight through. A one-wire box cannot inscribe |v_{jl}> without growing
% (tenkz spills the name), so the ket name sits at nw and the bra name at ne.
% The rank bound sum_j d_j <= C L^cbar is set under the right panel. Action
% runs right to left.
% Boundary signature: (phys:w)x3 | (phys:e)x3 on both sides of =sum_{j,l}
% (set as \sum\limits_{j,\ell}, J3).
% Label macros: none.
% Waiver: the split of the top two wires depends on j (X_j grows with j), so
% the left box is drawn in the coordinates of the j-th summand; the sum over j
% regroups the same three-wire signature. The summands have pairwise orthogonal
% ranges, and no bound on the internal complexity of v_{jl} is drawn or claimed.
% Related: peps_frame_hole_encoder_b and peps_hole_encoder draw the encoder
% K_{c,h}, whose K^dagger K is this projector with X_j -> D_j; no TNLean figure
% draws the projector sum itself.
\begin{tenkzequation}
    \tenkzkernel
    \begin{tenkz}[rows={wire,wire,wire}, cols=4, bonds=none, size=l]
        \tn[at=(1,2), name=cylP, skin=box, wide=2, wires=2, size=l,
            ports={180@1:physical, 180@2:physical, 0@1:physical, 0@2:physical}]{P}
        \tn[at=(1,1), name=cylWa, skin=none, ports={180:physical, 0:physical}]{}
        \tn[at=(2,1), name=cylWb, skin=none, ports={180:physical, 0:physical}]{}
        \tn[at=(3,1), name=cylWc, skin=none, ports={180:physical, 0:physical}]{}
        \tn[at=(1,4), name=cylEa, skin=none, ports={180:physical, 0:physical}]{}
        \tn[at=(2,4), name=cylEb, skin=none, ports={180:physical, 0:physical}]{}
        \tn[at=(3,4), name=cylEc, skin=none, ports={180:physical, 0:physical}]{}
        \tnwire{cylWa.0}{cylP.180@1}
        \tnwire{cylWb.0}{cylP.180@2}
        \tnwire{cylWc.0}{cylEc.180}
        \tnwire{cylP.0@1}{cylEa.180}
        \tnwire{cylP.0@2}{cylEb.180}
        \tnwire{cylEa.0}{open e}
        \tnwire{cylEb.0}{open e}
        \tnwire{cylEc.0}{open e}
        \tnmark[form=bracket, label pos=w]{(1,1) .. (2,1)}{$X(2u)$}
    \end{tenkz}
    $=\sum\limits_{j,\ell}$
    \begin{tenkz}[rows={wire,wire,wire}, cols=4, bonds=none, size=l]
        \tn[at=(1,2), name=cylKet, skin=box, size=l, label pos=nw,
            ports={180:physical}]{\ket{v_{j\ell}}}
        \tn[at=(1,3), name=cylBra, skin=box, size=l, label pos=ne,
            ports={0:physical}]{\bra{v_{j\ell}}}
        \tn[at=(1,1), name=cylRwa, skin=none, ports={180:physical, 0:physical}]{}
        \tn[at=(2,1), name=cylRwb, skin=none, ports={180:physical, 0:physical}]{}
        \tn[at=(3,1), name=cylRwc, skin=none, ports={180:physical, 0:physical}]{}
        \tn[at=(1,4), name=cylRea, skin=none, ports={180:physical, 0:physical}]{}
        \tn[at=(2,4), name=cylReb, skin=none, ports={180:physical, 0:physical}]{}
        \tn[at=(3,4), name=cylRec, skin=none, ports={180:physical, 0:physical}]{}
        \tnwire{open w}{cylRwa.180}
        \tnwire{cylRwa.0}{cylKet.180}
        \tnwire{cylBra.0}{cylRea.180}
        \tnwire[name=cylInIn]{cylRea.0}{open e}
        \tnwire{open w}{cylRwb.180}
        \tnwire{cylRwb.0}{cylReb.180}
        \tnwire[name=cylInPatch]{cylReb.0}{open e}
        \tnwire{open w}{cylRwc.180}
        \tnwire{cylRwc.0}{cylRec.180}
        \tnwire[name=cylInFar]{cylRec.0}{open e}
        \tnmark[form=label, label pos=e]{on cylInIn 1}{$X_j$}
        \tnmark[form=label, label pos=e]{on cylInPatch 1}{$X(2u)\setminus X_j$}
        \tnmark[form=label, label pos=e]{on cylInFar 1}{$\Lambda\setminus X(2u)$}
        \tnmark[form=label, label pos=s]{0.5 s of midway (3,2) and (3,3)}{$\sum_j d_j\le CL^{\bar c}$}
    \end{tenkz}
\end{tenkzequation}
